\documentclass[8pt,aspectratio=169]{beamer}
\usetheme{default}
\usecolortheme{dove}
\setbeamercovered{transparent}
\setbeamertemplate{navigation symbols}{} % remove navigation symbols

% \usepackage{fontspec} % font
% \setmainfont{Helvetica} % font

\usepackage{natbib} % citations + references
\bibliographystyle{apalike} % citations + references

\usepackage{mathrsfs} % equations

\title{Neuro-Dinner 2025}
\author[]{Gabriel Riegner}
\date{November 2025}

\begin{document}


% slide %
\begin{frame}{}
\textbf{Drift Diffusion Process}
\begin{align}
    Z_\tau = a z + \sum_{t=1}^\tau e_t, \quad e_t \sim \mathcal{N}(v\Delta \tau, \sigma^2 \Delta \tau)
\end{align}
$a > 0$ is the boundary height with decision boundaries at $\pm a$\\
$z \in [-1, +1]$ is the relative starting point \\
$v \in \mathbb{R}$ is the drift rate\\
$\sigma^2$ is the Gaussian noise variance\\
$\Delta \tau \to 0$ is a continuous drift diffusion process\\
$\tau$ is the diffusion time index with resolution $\Delta\tau$\\

\vfill
\textbf{Choice and Reaction Time for a Single Trial}
\begin{align}
    RT = \begin{cases}
        + \min\{\tau>0: Z_\tau \ge +a\} \cdot \Delta\tau + t_0\\
        - \min\{\tau>0: Z_\tau \le -a\} \cdot \Delta\tau + t_0
        \end{cases}
\end{align}
$t_0 \ge 0$ is the non-decision time\\
$\text{sign}(RT) \in \{-1, +1\}$ is the binary choice and $|RT|> 0$ is the reaction time\\
$+$ is upper boundary crossing, $-$ is lower boundary crossing

\vfill
\textbf{Choices and Reaction Times for Multiple Trials}
\begin{align}
    \{RT_i\}_{i=1}^n \sim f(RT;\theta = (a, t_0, v, z))
\end{align}

$i = 1, 2 \ldots,n$ is the trial index\\
$f$ is the joint choice and reaction time density determined by parameters $\theta$
\end{frame}

% slide %
\begin{frame}{}
\textbf{Probability Density Function} \citep{feller_introduction_1968, navarro_fast_2009}
\begin{align}
    f(RT ; \theta) = &\frac{4\pi}{a^2}\exp\left(\frac{\text{sign}(RT)avw-v^2(|RT|-t_0)}{2}\right)\\
    &\times\sum_{k=1}^\infty k\exp\left(-\frac{2k^2\pi^2 (|RT|-t_0)}{a^2}\right)\sin(k\pi w)\notag
\end{align}

$w = \frac{1-\text{sign}(RT)z}{2}$\\
$|RT| > t_0$


\vfill
\textbf{Log-Likelihood Function}
\begin{align}
    \ell_n (\theta ; \{RT_i\}_{i=1}^n) = \frac{1}{n}\sum_{i=1}^n \log f(RT_i ; \theta)
\end{align}

$\ell_n$ is the average log-likelihood over all trials

\vfill
\textbf{Maximum Likelihood Estimator}
\begin{align}
    \hat{\theta} = \underset{\theta}{\text{argmax}}\;\ell_n (\theta ; \{RT_i\}_{i=1}^n)
\end{align}

$\hat{\theta}$ is the parameter vector that maximizes the log-likelihood

\end{frame}


% slide %
\begin{frame}{}

\textbf{Hessian Matrix}
\begin{align}
        \widehat{\mathcal{H}} = -\frac{1}{n}\sum_{i=1}^n\frac{\partial^2}{\partial\theta\partial\theta^{\prime}}\log f(RT_i;\hat{\theta}) =-\frac{1}{n}\frac{\partial^2}{\partial\theta\partial\theta^{\prime}}\ell_n(\hat{\theta})
\end{align}
$\widehat{\mathcal{H}}$ is the sample Hessian, the matrix of second derivatives of the log-likelihood

\vfill
\textbf{Fisher Information Matrix}
\begin{align}
    \widehat{\mathscr{I}} = \frac{1}{n}\sum_{i=1}^n\left(\frac{\partial}{\partial\theta}\log f(RT_i;\hat{\theta})\right)\left(\frac{\partial}{\partial\theta}\log f(RT_i;\hat{\theta})\right)^{\prime} = \frac{1}{n} \sum_{i=1}^n \hat{S_i} \hat{S_i}^\prime
\end{align}
$\hat{S_i}$ is the score vector for trial $i$, the gradient of the log-likelihood\\
$\widehat{\mathscr{I}}$ is the sample Fisher information, the outer product of the scores

\end{frame}


\begin{frame}{}
\textbf{Sample Hessian}
\begin{align}
    \widehat V_\text{SH} = \widehat{\mathcal{H}}^{-1}
\end{align}
$\widehat V_\text{SH}$ is the covariance matrix estimated from the Hessian
\vfill

\textbf{Outer Product}
\begin{align}
    \widehat V_\text{OP} = \widehat{\mathscr{I}}^{-1}
\end{align}
$\widehat V_\text{OP}$ is the covariance matrix estimated from the Fisher information\\
\vfill

\textbf{Misspecification Robust} $\mathscr{I} \ne \mathcal{H}$\\
\begin{align}
    \widehat V_\text{MR} = \widehat{\mathcal{H}}_\theta^{-1} \widehat{\mathscr{I}}_\theta \widehat{\mathcal{H}}_\theta^{-1}
\end{align}
$\widehat V\text{MR}$ is the sandwich estimator, robust to model misspecification

\vfill
\textbf{Autocorrelation Robust}\\
\begin{align}
    \widehat V_\text{AR} = \widehat{\mathcal{H}}^{-1} \left(\frac{1}{n} \sum_{i,j=1}^n w_{|i-j|} \hat{S_i} \hat{S_j}^\prime\right) \widehat{\mathcal{H}}^{-1}
\end{align}
$\widehat V_\text{AR}$ is a covariance matrix estimator robust to autocorrelation\\
$w_{|i-j|}$ are weights that taper long-range autocorrelations

\end{frame}

% slide %
\begin{frame}{}

\textbf{Consistent Estimation} \citep[Theorem~10.8]{hansen_probability_2022}
\begin{align}
    \hat{\theta} \underset{p}{\to} \theta_0\;\text{as}\;n\to\infty
\end{align}

under \emph{Stationarity} and \emph{Ergodicity} \citep[Theorem~14.9]{hansen_econometrics_2022}\\
denotes convergence in probability of the estimator $\hat\theta$ to the true value $\theta_0$

\vfill
\textbf{Asymptotic Normality} \citep[Theorem~10.16]{hansen_probability_2022}
\begin{align}
    \sqrt{n}(\hat{\theta} - \theta_0) \underset{d}{\to} \mathcal{N}(0, V) \;\text{as}\;n\to\infty
\end{align}

under \emph{Strong Mixing} \citep[Theorem~14.15]{hansen_econometrics_2022}\\
denotes convergence in distribution to a Gaussian with mean 0 and covariance $V$
\end{frame}

\begin{frame}{}
    \begin{align}
        \hat{a}\; \bar{a}\\
        |\hat{v}|\; |\bar{v}|\\
        \hat{z}\; \bar{z}
    \end{align}

    \begin{align}
         \hat{f}(RT)\\
         f(RT; \theta^{(0)})\\
         f(RT; \hat\theta)
    \end{align}

    \begin{align}
         \ell_n (\theta=(a,v) ; \{RT_i\}_{i=1}^n)\\
         \ell_n (\theta=v ; \{RT_i\}_{i=1}^n)
    \end{align}

    \begin{align}
        a_i = a_{i-1} + \mathcal{N}(0, 0.02)
    \end{align}
\end{frame}

\begin{frame}{}
    \bibliography{docs/references}
\end{frame}

\end{document}
